\chapter{Implementación}
\label{cap:implementacion}

\section{Tecnologías y estructura}
El backend está implementado en Java 25 con Spring Boot 4.1.2 y Spring Cloud
2025.1.3. El cliente es una PWA Angular 22 e Ionic 9. Keycloak 26.8.1 gestiona
identidad y PostgreSQL 18.6 proporciona persistencia. La infraestructura local
incluye PostgreSQL, Keycloak y Mailpit. Estas versiones describen la línea de
implementación documentada en el repositorio.

\section{Flujo de emisión y recibo}
El servicio de acceso valida la credencial y consume un token de un solo uso.
La urna guarda el voto, el hash del recibo y el evento correspondiente en una
transacción. Si el almacenamiento falla, el token se revierte para permitir
un reintento controlado. El recibo acredita la presencia del voto en el
registro, pero no demuestra que la opción registrada coincida con la intención
del votante.

\section{Auditoría y escrutinio}
La urna encadena eventos mediante hashes y calcula una huella final. Antes del
escrutinio se verifica la cadena; una anomalía bloquea la publicación. El
recuento genera una huella del conjunto de votos y del resultado, que puede
recalcularse de forma independiente. Estas evidencias permiten comprobar
integridad y presencia, pero no constituyen verificabilidad criptográfica
extremo a extremo.

\section{Trazabilidad de requisitos}
La matriz RF-servicio-endpoint-pantalla y el estado de sus evidencias se
mantienen en \texttt{docs/ai/TRACEABILITY.md}. En ella constan, entre otros,
RF-08--RF-10 en voter-access, ballot-box y el flujo PWA; RF-12--RF-16 en
auditoría y resultados; y RF-17--RF-21 en notificaciones, privacidad,
trazabilidad y API de integración. La matriz conserva el detalle por endpoint
y evita duplicar una tabla extensa que se desactualizaría con facilidad.

\section{Material gráfico}
Las fuentes editables de secuencia y actividad se conservan en
\texttt{Figuras/diagramas/}. No se incorporan capturas nuevas en esta revisión:
las capturas representativas del cliente deben generarse en el entorno de
demostración y revisarse para asegurar que no contienen datos personales.
